\documentclass[11pt]{article}

\usepackage[margin=1in]{geometry}
\usepackage{amsmath,amssymb,amsthm}
\usepackage{booktabs}
\usepackage{longtable}
\usepackage{array}
\usepackage{graphicx}
\usepackage{hyperref}
\usepackage{authblk}
\usepackage{enumitem}
\usepackage[numbers,sort&compress]{natbib}

\hypersetup{colorlinks=true,linkcolor=blue,citecolor=blue,urlcolor=blue}

\newcommand{\J}[2]{J_{#1,#2}}
\newcommand{\ii}{\mathrm{i}}

\title{\textbf{Exact Traveling-Wave Solution Families of a Higher-Order \\
Gain--Loss Nonlinear Schr\"odinger-Type Equation via the iB-Function (BDKm) Ansatz: \\
A Combinatorial Solution Catalog and a Benchmark Design for \\
Physics-Informed Machine Learning}}

\author[1]{Author Name}
\author[1]{Co-Author Name}
\affil[1]{Department of Applied Mathematics / Photonics, Institution Name}

\date{\today}

\begin{document}

\maketitle

\begin{abstract}
We study a fourth-order nonlinear evolution equation that couples group-velocity
dispersion, third- and fourth-order dispersion, saturable Kerr-type self-phase
modulation, and gain/loss with finite gain-bandwidth, a model relevant to pulse
propagation in dispersive, dissipative, gain-saturated media such as doped fiber
amplifiers and mode-locked lasers. Using the traveling-wave reduction and the
implicit Bogning function (iB-function) ansatz $A(\xi)=a\,J_{n,m}(\xi)$ with
$J_{n,m}(\xi)=\sinh^m(\xi)/\cosh^n(\xi)$, we derive the fully expanded
nine-term-times-two (``Range'') algebraic equation obtained after coefficient
substitution, and we carry out a complete combinatorial classification of the
index pairs $(n,m)$ for which the two blocks of the Range equation can share
common $J$-indices. This classification yields a restricted field of 49 rational
half-integer pairs $(n,m)\in\{-3/2,-1,-1/2,0,1/2,1,3/2\}^2$, for which we tabulate
the fully expanded coefficients $C_1,\dots,C_9$ in closed form in terms of the
eight physical parameters of the model. Beyond the analytical contribution, we
propose and specify a concrete methodology for using this parametrized exact-solution
catalog as a benchmark and training-data generator for physics-informed neural
networks (PINNs), extending existing PINN studies of the nonlinear Schr\"odinger
hierarchy -- which typically validate against one or two closed-form solitons --
to a systematically generated, 49-branch family with known ground truth across the
full physical parameter set. We give the explicit forward-problem and
inverse-problem (parameter-discovery) loss functionals for this equation, work
through one representative branch ($n=m=1$) in full detail, and outline a
generalization-testing protocol (training on a subset of $(n,m)$ branches and
evaluating on held-out branches) that is not available with a single hand-picked
soliton solution. This paper is presented as a methodological framework and
worked derivation; the numerical PINN experiments it specifies are proposed future
work rather than results reported here.
\end{abstract}

\noindent\textbf{Keywords:} higher-order nonlinear Schr\"odinger equation; gain-saturated
dispersive media; traveling-wave ansatz; iB-functions; exact solutions; physics-informed
neural networks; benchmark datasets; parameter discovery

\tableofcontents

%======================================================================
\section{Introduction}
%======================================================================

Nonlinear evolution equations in the nonlinear Schr\"odinger (NLS) hierarchy govern
pulse propagation in a wide range of dispersive-nonlinear media, from optical fibers
to Bose--Einstein condensates and deep-water waves. When higher-order dispersion,
gain, loss, and saturable nonlinearity are all retained, the resulting model equation
becomes analytically intractable by inverse scattering in general, and one instead
seeks particular closed-form traveling-wave solutions through structured ans\"atze
(tanh-, sech-, Jacobi-elliptic-, and, more recently, ``implicit Bogning function''
or iB-function methods). These ansatz methods are algorithmic: a trial solution is
substituted into the governing equation, the result is expanded in a basis of
elementary functions, and the coefficients of each basis element are matched to
zero, producing a system of algebraic constraints on the physical and wave
parameters \citep{baldwin2004,hereman1990}.

Independently, over the last several years physics-informed neural networks (PINNs)
\citep{raissi2019} have become a standard tool for both the forward problem (learning
a solution field from the governing PDE and sparse data) and the inverse problem
(recovering unknown physical coefficients from field data) for equations in the NLS
hierarchy. PINNs have been applied to the standard NLS equation, the third-order
(Hirota) equation \citep{zhou2021}, higher-order NLS equations with self-steepening
and Raman terms \citep{fang2021}, equations with external potentials
\citep{pinnreview2025}, and adaptive-sampling variants that improve accuracy on
strongly localized structures such as rogue waves \citep{gravis2026}. A recurring
feature of this literature is that PINNs are trained and validated against a small
number of hand-selected closed-form solutions (one soliton, two solitons, a breather,
a rogue wave), each with its own fixed parameter values.

This paper makes two contributions. First, we present a complete and self-contained
derivation of the traveling-wave reduction of a fourth-order gain--loss NLS-type
equation using the iB-function ansatz $J_{n,m}(\xi)=\sinh^m(\xi)/\cosh^n(\xi)$, including
a systematic combinatorial classification of which index pairs $(n,m)$ can produce
internally consistent solutions once the ansatz is substituted and the resulting
expression is expanded over the saturable-nonlinearity denominator. This produces a
\emph{restricted field} of 49 admissible half-integer index pairs, each with fully
expanded, closed-form coefficients in terms of the eight physical parameters of the
model ($\beta_2,\beta_3,\beta_4,g_0,T_2,\alpha,\alpha_2,\gamma,\Gamma$) and the two
wave parameters $\lambda_1,\lambda_2$.

Second, we argue -- and specify in reproducible detail -- that this kind of
systematically generated, parametrized solution family is directly useful as a
benchmark and training-data generator for PINNs applied to higher-order gain-saturated
evolution equations, addressing a real gap in the existing PINN-for-NLSE literature:
existing studies validate against isolated solutions rather than families, so they
cannot test whether a learned solver or a parameter-discovery network generalizes
across a structured solution space. We give the explicit PINN forward and inverse loss
functionals for this equation (Section~\ref{sec:pinn}), work through the $(n,m)=(1,1)$
branch as a concrete worked example, and specify a generalization-testing protocol.
We do not report numerical PINN results in this paper; the machine-learning section is
a methodological proposal intended to make the benchmark immediately usable by
follow-up numerical work.

%======================================================================
\section{Governing Equation}
%======================================================================

We consider the fourth-order nonlinear evolution equation
\begin{equation}
\frac{\partial A}{\partial z}
+\frac{\ii}{2}\!\left(\beta_2+\ii g_0T_2^2\right)\frac{\partial^2 A}{\partial t^2}
+\frac{\beta_3}{6}\frac{\partial^3 A}{\partial t^3}
+\frac{\ii\beta_4}{24}\frac{\partial^4 A}{\partial t^4}
+\ii\!\left(\gamma+\ii\frac{\alpha_2}{2}\right)\frac{|A|^2A}{1+\Gamma|A|^2}
+\frac{1}{2}(g_0-\alpha)A=0,
\label{eq:governing}
\end{equation}
where $A(z,t)$ is the complex envelope, $z$ is the propagation coordinate, $t$ is the
retarded time, $\beta_2,\beta_3,\beta_4$ are the second-, third- and fourth-order
dispersion coefficients, $g_0$ is the small-signal gain, $T_2$ is the gain-medium
dipole relaxation time (setting the parabolic gain bandwidth), $\alpha$ is the linear
loss, $\alpha_2$ is a two-photon-type loss coefficient, $\gamma$ is the self-phase
modulation (Kerr) coefficient, and $\Gamma$ is the nonlinearity saturation parameter.
Equation~\eqref{eq:governing} reduces to the cubic NLS equation when
$\beta_3=\beta_4=g_0=\alpha=\alpha_2=\Gamma=0$, and to the Hirota/third-order NLS
equation when $\beta_4=g_0=\alpha_2=\Gamma=0$.

%======================================================================
\section{Traveling-Wave Reduction}
%======================================================================

We seek pure traveling-wave profiles by introducing
\begin{equation}
\xi=\lambda_1 z+\lambda_2 t, \qquad A(z,t)=A(\xi),
\label{eq:tw}
\end{equation}
with $\lambda_1,\lambda_2$ real propagation constants. The chain rule gives
\begin{align}
\frac{\partial A}{\partial z}&=\lambda_1\frac{dA}{d\xi}, &
\frac{\partial^2 A}{\partial t^2}&=\lambda_2^2\frac{d^2A}{d\xi^2}, \\
\frac{\partial^3 A}{\partial t^3}&=\lambda_2^3\frac{d^3A}{d\xi^3}, &
\frac{\partial^4 A}{\partial t^4}&=\lambda_2^4\frac{d^4A}{d\xi^4}.
\end{align}
Substituting into \eqref{eq:governing} yields the traveling-wave ODE
\begin{equation}
\lambda_1 A'
+\frac{\ii}{2}\!\left(\beta_2+\ii g_0T_2^2\right)\lambda_2^2 A''
+\frac{\beta_3}{6}\lambda_2^3A'''
+\frac{\ii\beta_4}{24}\lambda_2^4A^{(4)}
+\ii\!\left(\gamma+\ii\frac{\alpha_2}{2}\right)\frac{|A|^2A}{1+\Gamma|A|^2}
+\frac{1}{2}(g_0-\alpha)A=0,
\label{eq:tw-ode}
\end{equation}
where primes denote $d/d\xi$.

%======================================================================
\section{The iB-Function (BDKm) Ansatz}
%======================================================================

We represent the traveling-wave profile with the implicit Bogning function (iB-function)
ansatz,
\begin{equation}
A(\xi)=a\,J_{n,m}(\xi), \qquad
J_{n,m}(\xi)=\frac{\sinh^m(\xi)}{\cosh^n(\xi)},
\label{eq:ansatz}
\end{equation}
where $a$ is a (generally complex) amplitude and $n,m$ are exponents to be determined.
The ansatz obeys the closed differentiation rule
\begin{equation}
\frac{d}{d\xi}J_{n,m}=m\,J_{n-1,m-1}-n\,J_{n+1,m+1},
\label{eq:diffrule}
\end{equation}
which follows directly from $\dfrac{d}{d\xi}\sinh=\cosh$, $\dfrac{d}{d\xi}\cosh=\sinh$
and the quotient rule. Iterating \eqref{eq:diffrule} gives the derivatives of $A(\xi)$
used below.

\subsection{First and second derivatives}
\begin{align}
A' &= a\left[mJ_{n-1,m-1}-nJ_{n+1,m+1}\right], \\[2pt]
A'' &= a\Big[m(m-1)J_{n-2,m-2}-\big(m(n-1)+n(m+1)\big)J_{n,m}+n(n+1)J_{n+2,m+2}\Big].
\end{align}

\subsection{Third derivative}
\begin{equation}
A'''=a\left[A_1J_{n-3,m-3}+A_2J_{n-1,m-1}+A_3J_{n+1,m+1}+A_4J_{n+3,m+3}\right],
\end{equation}
with
\begin{align}
A_1&=m(m-1)(m-2)=m^3-3m^2+2m, \\
A_2&=-3m^2n+3m^2-2m, \\
A_3&=3mn^2+3n^2+2n, \\
A_4&=-n(n+1)(n+2)=-n^3-3n^2-2n.
\end{align}

\subsection{Fourth derivative}
\begin{align}
A^{(4)}=a\Big\{&A_1(m-3)J_{n-4,m-4}+\big[A_2(m-1)-A_1(n-3)\big]J_{n-2,m-2} \notag\\
&+\big[A_3(m+1)-A_2(n-1)\big]J_{n,m}
+\big[A_4(m+3)-A_3(n+1)\big]J_{n+2,m+2} \notag\\
&-A_4(n+3)J_{n+4,m+4}\Big\}.
\end{align}

%======================================================================
\section{Substitution and Coefficient Grouping}
%======================================================================

Using the power identities $|A|^2A=a^3J_{3n,3m}$ and $1+\Gamma|A|^2=1+\Gamma a^2J_{2n,2m}$
(taking $a$ real without loss of generality for the modulus), substituting the
derivatives above into \eqref{eq:tw-ode} and dividing by $a$ gives, before clearing the
saturation denominator,
\begin{equation}
\Big[\textstyle\sum_{k=1}^{9} C_k\,J_{n+r_k,\,m+r_k}\Big]
+\ii\!\left(\gamma+\ii\frac{\alpha_2}{2}\right)\frac{a^2J_{3n,3m}}{1+\Gamma a^2J_{2n,2m}}=0,
\label{eq:premultiplied}
\end{equation}
with offsets $r_k\in\{-4,-3,-2,-1,0,1,2,3,4\}$ for $k=1,\dots,9$ and the grouped
coefficients
\begin{align}
C_1&=\frac{\ii\beta_4}{24}\lambda_2^4\,m(m-1)(m-2)(m-3), \\
C_2&=\frac{\beta_3}{6}\lambda_2^3\,m(m-1)(m-2), \\
C_3&=\frac{\ii}{2}\!\left(\beta_2+\ii g_0T_2^2\right)\lambda_2^2m(m-1)
+\frac{\ii\beta_4}{24}\lambda_2^4\Big[(-3m^2n+3m^2-2m)(m-1)-m(m-1)(m-2)(n-3)\Big], \\
C_4&=\lambda_1 m+\frac{\beta_3}{6}\lambda_2^3\big(-3m^2n+3m^2-2m\big), \\
C_5&=\frac{1}{2}(g_0-\alpha)
-\frac{\ii}{2}\!\left(\beta_2+\ii g_0T_2^2\right)\lambda_2^2\big[m(n-1)+n(m+1)\big] \notag\\
&\quad+\frac{\ii\beta_4}{24}\lambda_2^4\Big[(3mn^2+3n^2+2n)(m+1)-(-3m^2n+3m^2-2m)(n-1)\Big], \\
C_6&=-\lambda_1 n+\frac{\beta_3}{6}\lambda_2^3\big(3mn^2+3n^2+2n\big), \\
C_7&=\frac{\ii}{2}\!\left(\beta_2+\ii g_0T_2^2\right)\lambda_2^2n(n+1)
+\frac{\ii\beta_4}{24}\lambda_2^4\Big[-n(n+1)(n+2)(m+3)-(3mn^2+3n^2+2n)(n+1)\Big], \\
C_8&=-\frac{\beta_3}{6}\lambda_2^3\,n(n+1)(n+2), \\
C_9&=\frac{\ii\beta_4}{24}\lambda_2^4\,n(n+1)(n+2)(n+3).
\end{align}

%======================================================================
\section{The Fully Expanded Range Equation}
%======================================================================

Multiplying \eqref{eq:premultiplied} by $1+\Gamma a^2J_{2n,2m}$ and using
$J_{2n,2m}\cdot J_{p,q}=J_{2n+p,2m+q}$ term by term gives the complete, fraction-free
19-term \emph{Range equation}:
\begin{align}
&C_1J_{n-4,m-4}+C_2J_{n-3,m-3}+C_3J_{n-2,m-2}+C_4J_{n-1,m-1}+C_5J_{n,m} \notag\\
&+C_6J_{n+1,m+1}+C_7J_{n+2,m+2}+C_8J_{n+3,m+3}+C_9J_{n+4,m+4} \notag\\
&+\Gamma a^2\Big[C_1J_{3n-4,3m-4}+C_2J_{3n-3,3m-3}+C_3J_{3n-2,3m-2}+C_4J_{3n-1,3m-1}+C_5J_{3n,3m} \notag\\
&\qquad+C_6J_{3n+1,3m+1}+C_7J_{3n+2,3m+2}+C_8J_{3n+3,3m+3}+C_9J_{3n+4,3m+4}\Big] \notag\\
&+\ii\!\left(\gamma+\ii\frac{\alpha_2}{2}\right)a^2J_{3n,3m}=0.
\label{eq:range}
\end{align}
For a solution to exist, every coefficient of every distinct $J_{p,q}$ index appearing
in \eqref{eq:range} must vanish (or the terms with matching indices must sum to zero).
Because the first block has indices $n+r$ ($r\in\{-4,\dots,4\}$) and the second block
has indices $3n+s$ ($s\in\{-4,\dots,4\}$), a term from each block coincides only for
special $(n,m)$; this motivates the combinatorial classification below.

%======================================================================
\section{Combinatorial Classification of Admissible $(n,m)$}
\label{sec:combinatorics}
%======================================================================

\subsection{Candidate values}

Let $r,s\in\{-4,-3,-2,-1,0,1,2,3,4\}$ index a term $J_{n+r,m+r}$ in the first block and
$J_{3n+s,3m+s}$ in the second block of \eqref{eq:range}. Equality of the first indices,
$n+r=3n+s$, requires
\begin{equation}
n=\frac{r-s}{2},
\label{eq:nsolve}
\end{equation}
and equality of the second indices gives the identical relation for $m$. Since a single
$J$-function pair must match simultaneously in both indices for the same $(r,s)$,
\begin{equation}
n=m=\frac{r-s}{2}.
\end{equation}
Table~\ref{tab:nvalues} tabulates $n=(r-s)/2$ over all $81$ ordered pairs $(r,s)$.

\begin{table}[htbp]
\centering
\caption{Values of $n=(r-s)/2$ obtained from $n+r=3n+s$, for $r,s\in\{-4,\dots,4\}$.}
\label{tab:nvalues}
\small
\begin{tabular}{c|ccccccccc}
\toprule
$r\backslash s$ & $-4$ & $-3$ & $-2$ & $-1$ & $0$ & $1$ & $2$ & $3$ & $4$ \\
\midrule
$-4$ & $0$    & $-\tfrac12$ & $-1$   & $-\tfrac32$ & $-2$ & $-\tfrac52$ & $-3$ & $-\tfrac72$ & $-4$ \\
$-3$ & $\tfrac12$ & $0$    & $-\tfrac12$ & $-1$   & $-\tfrac32$ & $-2$ & $-\tfrac52$ & $-3$ & $-\tfrac72$ \\
$-2$ & $1$    & $\tfrac12$ & $0$    & $-\tfrac12$ & $-1$   & $-\tfrac32$ & $-2$ & $-\tfrac52$ & $-3$ \\
$-1$ & $\tfrac32$ & $1$    & $\tfrac12$ & $0$    & $-\tfrac12$ & $-1$   & $-\tfrac32$ & $-2$ & $-\tfrac52$ \\
$0$  & $2$    & $\tfrac32$ & $1$    & $\tfrac12$ & $0$    & $-\tfrac12$ & $-1$   & $-\tfrac32$ & $-2$ \\
$1$  & $\tfrac52$ & $2$    & $\tfrac32$ & $1$    & $\tfrac12$ & $0$    & $-\tfrac12$ & $-1$   & $-\tfrac32$ \\
$2$  & $3$    & $\tfrac52$ & $2$    & $\tfrac32$ & $1$    & $\tfrac12$ & $0$    & $-\tfrac12$ & $-1$ \\
$3$  & $\tfrac72$ & $3$    & $\tfrac52$ & $2$    & $\tfrac32$ & $1$    & $\tfrac12$ & $0$    & $-\tfrac12$ \\
$4$  & $4$    & $\tfrac72$ & $3$    & $\tfrac52$ & $2$    & $\tfrac32$ & $1$    & $\tfrac12$ & $0$ \\
\bottomrule
\end{tabular}
\end{table}

Although there are 81 ordered comparisons, they collapse to only 17 distinct candidate
pairs $(n,m)=(q,q)$ with
\begin{equation}
q\in\left\{-4,-\tfrac72,-3,-\tfrac52,-2,-\tfrac32,-1,-\tfrac12,0,\tfrac12,1,\tfrac32,2,\tfrac52,3,\tfrac72,4\right\}.
\end{equation}
Restricting to integer $(n,m)$ leaves $q=-4,\dots,4$; further restricting to nonnegative
integers leaves $q=0,1,2,3,4$.

\subsection{Combinatorial frequency of each candidate}

Assigning each of the 81 ordered comparisons $(r,s)$ equal weight $1/81$, the number of
comparisons yielding a given difference $d=r-s=2q$ is $N_d=9-|d|$, so
\begin{equation}
N(q)=9-2|q|, \qquad P(q,q)=\frac{N(q)}{81}=\frac{9-2|q|}{81},
\qquad q=-4,-\tfrac72,\dots,\tfrac72,4.
\label{eq:prob}
\end{equation}
We stress that $P(q,q)$ is a purely combinatorial frequency over index comparisons, not a
physical probability of solution existence; it is a convenient measure of how many
index-matching routes support a given candidate pair. Table~\ref{tab:freq} lists the full
distribution, which sums to unity by construction.

\begin{table}[htbp]
\centering
\caption{Combinatorial frequency of each candidate pair $(n,m)=(q,q)$.}
\label{tab:freq}
\begin{tabular}{cccc}
\toprule
$(n,m)$ & Appearances & Probability & Percentage \\
\midrule
$(-4,-4)$ & 1 & $1/81$ & 1.23\% \\
$(-7/2,-7/2)$ & 2 & $2/81$ & 2.47\% \\
$(-3,-3)$ & 3 & $3/81$ & 3.70\% \\
$(-5/2,-5/2)$ & 4 & $4/81$ & 4.94\% \\
$(-2,-2)$ & 5 & $5/81$ & 6.17\% \\
$(-3/2,-3/2)$ & 6 & $6/81$ & 7.41\% \\
$(-1,-1)$ & 7 & $7/81$ & 8.64\% \\
$(-1/2,-1/2)$ & 8 & $8/81$ & 9.88\% \\
$(0,0)$ & 9 & $9/81$ & 11.11\% \\
$(1/2,1/2)$ & 8 & $8/81$ & 9.88\% \\
$(1,1)$ & 7 & $7/81$ & 8.64\% \\
$(3/2,3/2)$ & 6 & $6/81$ & 7.41\% \\
$(2,2)$ & 5 & $5/81$ & 6.17\% \\
$(5/2,5/2)$ & 4 & $4/81$ & 4.94\% \\
$(3,3)$ & 3 & $3/81$ & 3.70\% \\
$(7/2,7/2)$ & 2 & $2/81$ & 2.47\% \\
$(4,4)$ & 1 & $1/81$ & 1.23\% \\
\bottomrule
\end{tabular}
\end{table}

The most combinatorially supported candidate is $(0,0)$, corresponding to the trivial
constant $J_{0,0}\equiv1$; the extreme pairs $(\pm4,\pm4)$ are supported by a single
comparison route each.

\subsection{The restricted field}

The full extended field formed from the 17 candidate values in each coordinate contains
$17^2=289$ pairs. We restrict attention to the field generated by the six most
combinatorially supported nontrivial candidates together with $(0,0)$,
\begin{equation}
\mathcal{F}_R=\left\{-\tfrac32,-1,-\tfrac12,0,\tfrac12,1,\tfrac32\right\}^2,
\end{equation}
which contains $49$ ordered pairs $(n,m)$. For every pair in $\mathcal{F}_R$ the coefficients
$C_1,\dots,C_9$ can be evaluated by exact substitution of $n,m$ into the expressions of
Section~5. The full 49-branch table is given in Appendix~\ref{app:tables}.

%======================================================================
\section{A Worked Branch: $(n,m)=(1,1)$}
\label{sec:worked}
%======================================================================

To make the general procedure concrete, we work through the branch $(n,m)=(1,1)$, for
which $J_{1,1}(\xi)=\tanh(\xi)$, so the ansatz is the familiar dark/kink profile
$A(\xi)=a\tanh(\xi)$. Evaluating the coefficient expressions of Section~5 at $n=m=1$
gives (Appendix~\ref{app:tables}, Table~\ref{tab:n1-0113})
\begin{align}
C_1&=0, & C_2&=0, & C_3&=0, \\
C_4&=\lambda_1-\frac{\beta_3}{3}\lambda_2^3, &
C_5&=\frac{g_0-\alpha}{2}-\ii\!\left(\beta_2+\ii\frac{g_0T_2^2}{2}\right)\lambda_2^2+\frac{2\ii\beta_4}{3}\lambda_2^4, \\
C_6&=-\lambda_1+\frac{4\beta_3}{3}\lambda_2^3, &
C_7&=\ii\!\left(\beta_2+\ii\frac{g_0T_2^2}{2}\right)\lambda_2^2-\frac{5\ii\beta_4}{3}\lambda_2^4, \\
C_8&=-\beta_3\lambda_2^3, & C_9&=\ii\beta_4\lambda_2^4.
\end{align}
The vanishing of $C_1,C_2,C_3$ is not a coincidence: for $n=m=1$ the offset indices
$n-4,n-3,n-2$ correspond to $J_{-3,-3},J_{-2,-2},J_{-1,-1}$, and the prefactors $A_1(m-3)$,
$A_2(m-1)-A_1(n-3)$, and the second-derivative bracket all contain a factor of
$(m-1)$ or $(m-1)(m-2)\cdots$ that vanishes identically at $m=1$. This is a generic
feature of integer branches with $n=m$: the ``deep negative offset'' terms of the fourth
derivative collapse automatically, leaving a shorter, more tractable coefficient system.
The remaining Range equation for this branch, after also imposing the matching of
$J_{3,3}=\tanh^3(\xi)$ terms coming from the saturable-nonlinearity expansion and the
self-phase-modulation term, is
\begin{equation}
C_4J_{0,0}+C_5J_{1,1}+C_6J_{2,2}+C_7J_{2,2}
+\Gamma a^2\big[\cdots\big]_{J_{3,3},\,J_{5,5},\dots}
+\ii\!\left(\gamma+\ii\frac{\alpha_2}{2}\right)a^2J_{3,3}=0,
\end{equation}
and setting the coefficient of each independent power of $\tanh(\xi)$ to zero yields an
algebraic system relating $a,\lambda_1$ to $\beta_2,\beta_3,\beta_4,g_0,T_2,\alpha,\alpha_2,\gamma,\Gamma$
(most directly: $C_4=0$ fixes $\lambda_1=\beta_3\lambda_2^3/3$ once the cubic-dispersion
branch is selected, and the remaining balance between $C_5+C_6+C_7$ and the
nonlinear/saturation terms fixes $a^2$ in terms of the gain-loss and dispersion
parameters). We do not carry every branch of $\mathcal{F}_R$ to this level of explicit
closure in this paper; the point of Table~\ref{tab:freq}--\ref{tab:n32-3-32} is to make the
coefficient system for \emph{any} of the 49 branches immediately available for such a
reduction, whether performed by hand or automatically with a computer-algebra or
symbolic-regression back end (Section~\ref{sec:pinn}).

%======================================================================
\section{Application: A Benchmark Design for Physics-Informed Neural Networks}
\label{sec:pinn}
%======================================================================

\subsection{Motivation}

PINNs \citep{raissi2019} approximate the field $A(z,t)=u(z,t)+\ii v(z,t)$ by a neural
network $\hat A_\theta$ and train it by minimizing a composite loss built from (i) a
data-fit term against known values (initial/boundary conditions, or sparse
measurements) and (ii) a physics-residual term obtained by substituting $\hat A_\theta$
into the governing PDE and evaluating the residual on collocation points via automatic
differentiation. This scheme has been applied successfully to the NLS equation, the
Hirota (third-order NLS) equation \citep{zhou2021}, and higher-order NLS equations with
self-steepening and Raman terms \citep{fang2021}, in each case validated against one or
a handful of hand-picked exact solutions with fixed coefficients. A structural
limitation of this approach is that it cannot easily test \emph{generalization}: whether
a network (or a class of network architectures/training schedules) that performs well on
one soliton also performs well across a structured family of solutions with different
exponent structure and different physical-parameter regimes.

The restricted field $\mathcal{F}_R$ derived in Section~\ref{sec:combinatorics}
directly addresses this limitation for the higher-order gain-saturated equation
\eqref{eq:governing}: it provides 49 closed-form solution branches, each an exact
solution of \eqref{eq:governing} once the associated algebraic constraint system (as
illustrated in Section~\ref{sec:worked}) is imposed on
$(a,\lambda_1,\lambda_2,\beta_2,\beta_3,\beta_4,g_0,T_2,\alpha,\alpha_2,\gamma,\Gamma)$.
Because the branches differ systematically in $(n,m)$ -- and hence in qualitative
solution shape (bright, dark, singular, or mixed hyperbolic profiles) -- they form a
natural train/held-out split for generalization testing that a single soliton solution
cannot provide.

\subsection{Forward-problem loss functional}

Writing $A=u+\ii v$ and the residual $F=F_u+\ii F_v$ obtained from
\eqref{eq:governing},
\begin{align}
F_u&:=u_z-\tfrac12\!\left(\beta_2 v_{tt}+g_0T_2^2u_{tt}\right)-\tfrac{\beta_3}{6}v_{ttt}-\tfrac{\beta_4}{24}u_{tttt}
-\left(\gamma+\ii\tfrac{\alpha_2}{2}\right)\!\Re\!\left[\ii\frac{(u^2+v^2)(u+\ii v)}{1+\Gamma(u^2+v^2)}\right]
+\tfrac12(g_0-\alpha)u, \\
F_v&:=v_z+\tfrac12\!\left(\beta_2 u_{tt}-g_0T_2^2v_{tt}\right)+\tfrac{\beta_3}{6}u_{ttt}-\tfrac{\beta_4}{24}v_{tttt}
-\left(\gamma+\ii\tfrac{\alpha_2}{2}\right)\!\Im\!\left[\ii\frac{(u^2+v^2)(u+\ii v)}{1+\Gamma(u^2+v^2)}\right]
+\tfrac12(g_0-\alpha)v,
\end{align}
the forward-problem training loss for a given branch $(n,m)\in\mathcal{F}_R$, with its
exact solution $A_{n,m}(\xi;\,\text{params})$ supplying initial/boundary data, is
\begin{align}
\mathcal{L}_{\text{fwd}}=\;&\frac{1}{N_I}\sum_{j=1}^{N_I}\big|\hat A_\theta(z_0,t_I^j)-A_{n,m}(z_0,t_I^j)\big|^2
+\frac{1}{N_B}\sum_{j=1}^{N_B}\big|\hat A_\theta(z,\pm L_j)-A_{n,m}(z,\pm L_j)\big|^2 \notag\\
&+\frac{1}{N_S}\sum_{j=1}^{N_S}\big(F_u^{2}+F_v^{2}\big)\Big|_{(z^j_S,t^j_S)},
\end{align}
following the standard Latin-hypercube collocation strategy used throughout this
literature \citep{raissi2019,zhou2021,fang2021}.

\subsection{Inverse-problem (parameter-discovery) loss functional}

For parameter discovery, $\beta_2,\beta_3,\beta_4,g_0,T_2,\alpha,\alpha_2,\gamma,\Gamma$
(or any chosen subset) are promoted to trainable scalars alongside the network weights,
and the loss becomes
\begin{align}
\mathcal{L}_{\text{inv}}=\;&\frac{1}{N_p}\sum_{j=1}^{N_p}\big|\hat A_\theta(z^j,t^j)-A_{n,m}^j\big|^2
+\frac{1}{N_p}\sum_{j=1}^{N_p}\big(F_u^2+F_v^2\big)\Big|_{(z^j,t^j)},
\end{align}
with $\{z^j,t^j,A_{n,m}^j\}$ sampled from one exact branch. This is a direct
generalization of the two- and four-parameter discovery experiments of
\citet{zhou2021} and \citet{fang2021} to a model with up to nine independent physical
coefficients; the restricted field $\mathcal{F}_R$ makes it possible to test parameter
recovery separately on bright-type branches (e.g.\ small $|n|,|m|$) and singular/kink-type
branches (e.g.\ $n=m=1$), which is not possible with a single soliton family.

\subsection{Proposed generalization-testing protocol}

We propose the following protocol as future numerical work built directly on the
catalog in Appendix~\ref{app:tables}:
\begin{enumerate}[label=(\roman*)]
\item \textbf{Branch generation.} For each $(n,m)\in\mathcal{F}_R$, solve the algebraic
closure conditions illustrated in Section~\ref{sec:worked} (symbolically, e.g.\ with a
computer-algebra system) to obtain the admissible parameter manifold and one or more
representative parameter sets; sample $\{z^j,t^j,A^j\}$ from each resulting exact solution.
\item \textbf{Held-out split.} Partition $\mathcal{F}_R$ into a training subset (e.g.\
$|n|,|m|\le 1$) and a held-out subset (e.g.\ $|n|$ or $|m|=3/2$), train a single PINN
architecture on the training subset only, and evaluate the relative $L^2$ error on the
held-out branches.
\item \textbf{Parameter-recovery robustness.} Repeat the inverse-problem experiment of
\citet{zhou2021,fang2021} independently on each branch and report how recovery error for
$(\beta_2,\beta_3,\beta_4,g_0,T_2,\alpha,\alpha_2,\gamma,\Gamma)$ varies with branch
type, under both clean and noise-perturbed data (as in \citet{zhou2021}).
\item \textbf{Ablation on saturation.} Because $\Gamma$ enters only through the
denominator $1+\Gamma|A|^2$, compare PINN performance with $\Gamma=0$ (unsaturated,
standard higher-order NLS) against $\Gamma\neq0$ branches, to quantify whether the
saturable nonlinearity increases the forward- or inverse-problem difficulty.
\end{enumerate}
This protocol, and the underlying 49-branch catalog, is the concrete deliverable this
paper contributes to future PINN work on higher-order gain-saturated evolution
equations; we emphasize again that we report the framework and worked derivation only,
not executed numerical experiments.

%======================================================================
\section{Discussion}
%======================================================================

The combinatorial classification of Section~\ref{sec:combinatorics} is, to our
knowledge, not typically made explicit in the ansatz-method literature, where
authors usually state the admissible $(n,m)$ (or equivalent tanh/sech exponents) without
enumerating the full comparison space or reporting a frequency measure over it. Making
this step explicit and reproducible is what allows the resulting solution catalog to be
generated systematically (Appendix~\ref{app:tables}) rather than one branch at a time,
which is what makes it usable as a benchmark generator rather than a single validation
example. Existing symbolic-computation packages for the tanh/sech method
\citep{baldwin2004} automate a closely related search over exponents for a narrower
ansatz family; extending that kind of automation to the iB-function family
$J_{n,m}$ with fractional exponents, and coupling it directly to PINN-benchmark
generation as specified in Section~\ref{sec:pinn}, is a natural direction for follow-up
software work.

We also note two limitations. First, not every branch in $\mathcal{F}_R$ will admit a
non-degenerate closure (some coefficient systems may be over- or under-determined,
forcing $a=0$ or leaving free parameters unconstrained); a systematic audit of which of
the 49 branches yield genuine, non-trivial exact solutions is needed before the full
catalog can be used as-is as a benchmark suite. Second, the combinatorial frequency
$P(q,q)$ of Eq.~\eqref{eq:prob} measures index-matching multiplicity, not solution
existence or physical relevance, and should not be over-interpreted as a likelihood that
a given branch yields a physically meaningful pulse.

%======================================================================
\section{Conclusion}
%======================================================================

We have derived the traveling-wave reduction of a fourth-order gain-saturated nonlinear
Schr\"odinger-type equation under the iB-function (BDKm) ansatz, obtained the fully
expanded Range equation, and carried out a complete combinatorial classification of the
admissible index pairs $(n,m)$, yielding a restricted field of 49 branches with fully
expanded closed-form coefficients in terms of all eight physical parameters of the
model. We have shown, through the explicit $(n,m)=(1,1)$ example, how each branch reduces
to an algebraic closure system fixing the wave parameters in terms of the physical
ones. Building on this catalog, we specified a concrete benchmark design for
physics-informed neural networks -- forward-problem and inverse-problem loss functionals
matching the full nine-parameter model, and a generalization-testing protocol exploiting
the branch structure of $\mathcal{F}_R$ -- that addresses a specific limitation of
existing PINN-for-NLSE-hierarchy studies, namely their reliance on isolated,
hand-picked solutions rather than structured families. Executing this protocol
numerically, and auditing the non-degeneracy of all 49 branches, are the natural next
steps.

%======================================================================
\appendix
\section{Fully Expanded Coefficient Tables for the Restricted Field}
\label{app:tables}
%======================================================================

For every $(n,m)\in\mathcal{F}_R=\{-3/2,-1,-1/2,0,1/2,1,3/2\}^2$ we report the exact,
simplified values of $C_1,\dots,C_9$ from Section~5, grouped by the index $n$.

\subsection{$n=-3/2$}

\begin{table}[htbp]
\centering\small
\caption{Coefficients for $n=-3/2$, $m\in\{-3/2,-1,-1/2\}$.}
\begin{tabular}{lccc}
\toprule
Coeff. & $(-3/2,-3/2)$ & $(-3/2,-1)$ & $(-3/2,-1/2)$ \\
\midrule
$C_1$ & $\tfrac{315\ii\beta_4}{128}\lambda_2^4$ & $\ii\beta_4\lambda_2^4$ & $\tfrac{35\ii\beta_4}{128}\lambda_2^4$ \\
$C_2$ & $-\tfrac{35\beta_3}{16}\lambda_2^3$ & $-\beta_3\lambda_2^3$ & $-\tfrac{5\beta_3}{16}\lambda_2^3$ \\
$C_3$ & $\tfrac{15\ii}{8}\Xi\lambda_2^2-\tfrac{145\ii\beta_4}{32}\lambda_2^4$ & $\ii\Xi\lambda_2^2-\tfrac{23\ii\beta_4}{12}\lambda_2^4$ & $\tfrac{3\ii}{8}\Xi\lambda_2^2-\tfrac{17\ii\beta_4}{32}\lambda_2^4$ \\
$C_4$ & $-\tfrac32\lambda_1+\tfrac{53\beta_3}{16}\lambda_2^3$ & $-\lambda_1+\tfrac{19\beta_3}{12}\lambda_2^3$ & $-\tfrac12\lambda_1+\tfrac{23\beta_3}{48}\lambda_2^3$ \\
$C_5$ & $\tfrac{g_0-\alpha}{2}-\tfrac{9\ii}{4}\Xi\lambda_2^2+\tfrac{141\ii\beta_4}{64}\lambda_2^4$ & $\tfrac{g_0-\alpha}{2}-\tfrac{5\ii}{4}\Xi\lambda_2^2+\tfrac{95\ii\beta_4}{96}\lambda_2^4$ & $\tfrac{g_0-\alpha}{2}-\tfrac{\ii}{4}\Xi\lambda_2^2+\tfrac{59\ii\beta_4}{192}\lambda_2^4$ \\
$C_6$ & $\tfrac32\lambda_1-\tfrac{17\beta_3}{16}\lambda_2^3$ & $\tfrac32\lambda_1-\tfrac{\beta_3}{2}\lambda_2^3$ & $\tfrac32\lambda_1+\tfrac{\beta_3}{16}\lambda_2^3$ \\
$C_7$ & $\tfrac{3\ii}{8}\Xi\lambda_2^2-\tfrac{5\ii\beta_4}{32}\lambda_2^4$ & $\tfrac{3\ii}{8}\Xi\lambda_2^2-\tfrac{3\ii\beta_4}{32}\lambda_2^4$ & $\tfrac{3\ii}{8}\Xi\lambda_2^2-\tfrac{\ii\beta_4}{32}\lambda_2^4$ \\
$C_8$ & $-\tfrac{\beta_3}{16}\lambda_2^3$ & $-\tfrac{\beta_3}{16}\lambda_2^3$ & $-\tfrac{\beta_3}{16}\lambda_2^3$ \\
$C_9$ & $\tfrac{3\ii\beta_4}{128}\lambda_2^4$ & $\tfrac{3\ii\beta_4}{128}\lambda_2^4$ & $\tfrac{3\ii\beta_4}{128}\lambda_2^4$ \\
\bottomrule
\end{tabular}
\vspace{2pt}
\footnotesize{$\Xi:=\beta_2+\ii g_0T_2^2$ throughout the appendix.}
\end{table}

\begin{table}[htbp]
\centering\small
\caption{Coefficients for $n=-3/2$, $m\in\{0,1/2,1,3/2\}$.}
\begin{tabular}{lcccc}
\toprule
Coeff. & $(-3/2,0)$ & $(-3/2,1/2)$ & $(-3/2,1)$ & $(-3/2,3/2)$ \\
\midrule
$C_1$ & $0$ & $-\tfrac{5\ii\beta_4}{128}\lambda_2^4$ & $0$ & $\tfrac{3\ii\beta_4}{128}\lambda_2^4$ \\
$C_2$ & $0$ & $\tfrac{\beta_3}{16}\lambda_2^3$ & $0$ & $-\tfrac{\beta_3}{16}\lambda_2^3$ \\
$C_3$ & $0$ & $-\tfrac{\ii}{8}\Xi\lambda_2^2+\tfrac{5\ii\beta_4}{96}\lambda_2^4$ & $0$ & $\tfrac{3\ii}{8}\Xi\lambda_2^2+\tfrac{7\ii\beta_4}{32}\lambda_2^4$ \\
$C_4$ & $0$ & $\tfrac12\lambda_1+\tfrac{7\beta_3}{48}\lambda_2^3$ & $\lambda_1+\tfrac{11\beta_3}{12}\lambda_2^3$ & $\tfrac32\lambda_1+\tfrac{37\beta_3}{16}\lambda_2^3$ \\
$C_5$ & $\tfrac{g_0-\alpha}{2}+\tfrac{3\ii}{4}\Xi\lambda_2^2+\tfrac{5\ii\beta_4}{32}\lambda_2^4$ & $\tfrac{g_0-\alpha}{2}+\tfrac{7\ii}{4}\Xi\lambda_2^2+\tfrac{103\ii\beta_4}{192}\lambda_2^4$ & $\tfrac{g_0-\alpha}{2}+\tfrac{11\ii}{4}\Xi\lambda_2^2+\tfrac{139\ii\beta_4}{96}\lambda_2^4$ & $\tfrac{g_0-\alpha}{2}+\tfrac{15\ii}{4}\Xi\lambda_2^2+\tfrac{185\ii\beta_4}{64}\lambda_2^4$ \\
$C_6$ & $\tfrac32\lambda_1+\tfrac{5\beta_3}{8}\lambda_2^3$ & $\tfrac32\lambda_1+\tfrac{19\beta_3}{16}\lambda_2^3$ & $\tfrac32\lambda_1+\tfrac{7\beta_3}{4}\lambda_2^3$ & $\tfrac32\lambda_1+\tfrac{37\beta_3}{16}\lambda_2^3$ \\
$C_7$ & $\tfrac{3\ii}{8}\Xi\lambda_2^2+\tfrac{\ii\beta_4}{32}\lambda_2^4$ & $\tfrac{3\ii}{8}\Xi\lambda_2^2+\tfrac{3\ii\beta_4}{32}\lambda_2^4$ & $\tfrac{3\ii}{8}\Xi\lambda_2^2+\tfrac{5\ii\beta_4}{32}\lambda_2^4$ & $\tfrac{3\ii}{8}\Xi\lambda_2^2+\tfrac{7\ii\beta_4}{32}\lambda_2^4$ \\
$C_8$ & $-\tfrac{\beta_3}{16}\lambda_2^3$ & $-\tfrac{\beta_3}{16}\lambda_2^3$ & $-\tfrac{\beta_3}{16}\lambda_2^3$ & $-\tfrac{\beta_3}{16}\lambda_2^3$ \\
$C_9$ & $\tfrac{3\ii\beta_4}{128}\lambda_2^4$ & $\tfrac{3\ii\beta_4}{128}\lambda_2^4$ & $\tfrac{3\ii\beta_4}{128}\lambda_2^4$ & $\tfrac{3\ii\beta_4}{128}\lambda_2^4$ \\
\bottomrule
\end{tabular}
\end{table}

\subsection{$n=-1$}

\begin{table}[htbp]
\centering\small
\caption{Coefficients for $n=-1$, $m\in\{-3/2,-1,-1/2\}$.}
\begin{tabular}{lccc}
\toprule
Coeff. & $(-1,-3/2)$ & $(-1,-1)$ & $(-1,-1/2)$ \\
\midrule
$C_1$ & $\tfrac{315\ii\beta_4}{128}\lambda_2^4$ & $\ii\beta_4\lambda_2^4$ & $\tfrac{35\ii\beta_4}{128}\lambda_2^4$ \\
$C_2$ & $-\tfrac{35\beta_3}{16}\lambda_2^3$ & $-\beta_3\lambda_2^3$ & $-\tfrac{5\beta_3}{16}\lambda_2^3$ \\
$C_3$ & $\tfrac{15\ii}{8}\Xi\lambda_2^2-\tfrac{125\ii\beta_4}{32}\lambda_2^4$ & $\ii\Xi\lambda_2^2-\tfrac{5\ii\beta_4}{3}\lambda_2^4$ & $\tfrac{3\ii}{8}\Xi\lambda_2^2-\tfrac{15\ii\beta_4}{32}\lambda_2^4$ \\
$C_4$ & $-\tfrac32\lambda_1+\tfrac{11\beta_3}{4}\lambda_2^3$ & $-\lambda_1+\tfrac{4\beta_3}{3}\lambda_2^3$ & $-\tfrac12\lambda_1+\tfrac{5\beta_3}{12}\lambda_2^3$ \\
$C_5$ & $\tfrac{g_0-\alpha}{2}-\tfrac{7\ii}{4}\Xi\lambda_2^2+\tfrac{139\ii\beta_4}{96}\lambda_2^4$ & $\tfrac{g_0-\alpha}{2}-\ii\Xi\lambda_2^2+\tfrac{2\ii\beta_4}{3}\lambda_2^4$ & $\tfrac{g_0-\alpha}{2}-\tfrac{\ii}{4}\Xi\lambda_2^2+\tfrac{19\ii\beta_4}{96}\lambda_2^4$ \\
$C_6$ & $\lambda_1-\tfrac{7\beta_3}{12}\lambda_2^3$ & $\lambda_1-\tfrac{\beta_3}{3}\lambda_2^3$ & $\lambda_1-\tfrac{\beta_3}{12}\lambda_2^3$ \\
$C_7$ & $0$ & $0$ & $0$ \\
$C_8$ & $0$ & $0$ & $0$ \\
$C_9$ & $0$ & $0$ & $0$ \\
\bottomrule
\end{tabular}
\end{table}

\begin{table}[htbp]
\centering\small
\caption{Coefficients for $n=-1$, $m\in\{0,1/2,1,3/2\}$.}
\label{tab:nneg1-0}
\begin{tabular}{lcccc}
\toprule
Coeff. & $(-1,0)$ & $(-1,1/2)$ & $(-1,1)$ & $(-1,3/2)$ \\
\midrule
$C_1$ & $0$ & $-\tfrac{5\ii\beta_4}{128}\lambda_2^4$ & $0$ & $\tfrac{3\ii\beta_4}{128}\lambda_2^4$ \\
$C_2$ & $0$ & $\tfrac{\beta_3}{16}\lambda_2^3$ & $0$ & $-\tfrac{\beta_3}{16}\lambda_2^3$ \\
$C_3$ & $0$ & $-\tfrac{\ii}{8}\Xi\lambda_2^2+\tfrac{5\ii\beta_4}{96}\lambda_2^4$ & $0$ & $\tfrac{3\ii}{8}\Xi\lambda_2^2+\tfrac{5\ii\beta_4}{32}\lambda_2^4$ \\
$C_4$ & $0$ & $\tfrac12\lambda_1+\tfrac{\beta_3}{12}\lambda_2^3$ & $\lambda_1+\tfrac{2\beta_3}{3}\lambda_2^3$ & $\tfrac32\lambda_1+\tfrac{7\beta_3}{4}\lambda_2^3$ \\
$C_5$ & $\tfrac{g_0-\alpha}{2}+\tfrac{\ii}{2}\Xi\lambda_2^2+\tfrac{\ii\beta_4}{24}\lambda_2^4$ & $\tfrac{g_0-\alpha}{2}+\tfrac{5\ii}{4}\Xi\lambda_2^2+\tfrac{19\ii\beta_4}{96}\lambda_2^4$ & $\tfrac{g_0-\alpha}{2}+2\ii\Xi\lambda_2^2+\tfrac{2\ii\beta_4}{3}\lambda_2^4$ & $\tfrac{g_0-\alpha}{2}+\tfrac{11\ii}{4}\Xi\lambda_2^2+\tfrac{139\ii\beta_4}{96}\lambda_2^4$ \\
$C_6$ & $\lambda_1+\tfrac{\beta_3}{6}\lambda_2^3$ & $\lambda_1+\tfrac{5\beta_3}{12}\lambda_2^3$ & $\lambda_1+\tfrac{2\beta_3}{3}\lambda_2^3$ & $\lambda_1+\tfrac{11\beta_3}{12}\lambda_2^3$ \\
$C_7$ & $0$ & $0$ & $0$ & $0$ \\
$C_8$ & $0$ & $0$ & $0$ & $0$ \\
$C_9$ & $0$ & $0$ & $0$ & $0$ \\
\bottomrule
\end{tabular}
\end{table}

\subsection{$n=-1/2$}

\begin{table}[htbp]
\centering\small
\caption{Coefficients for $n=-1/2$, $m\in\{-3/2,-1,-1/2\}$.}
\begin{tabular}{lccc}
\toprule
Coeff. & $(-1/2,-3/2)$ & $(-1/2,-1)$ & $(-1/2,-1/2)$ \\
\midrule
$C_1$ & $\tfrac{315\ii\beta_4}{128}\lambda_2^4$ & $\ii\beta_4\lambda_2^4$ & $\tfrac{35\ii\beta_4}{128}\lambda_2^4$ \\
$C_2$ & $-\tfrac{35\beta_3}{16}\lambda_2^3$ & $-\beta_3\lambda_2^3$ & $-\tfrac{5\beta_3}{16}\lambda_2^3$ \\
$C_3$ & $\tfrac{15\ii}{8}\Xi\lambda_2^2-\tfrac{105\ii\beta_4}{32}\lambda_2^4$ & $\ii\Xi\lambda_2^2-\tfrac{17\ii\beta_4}{12}\lambda_2^4$ & $\tfrac{3\ii}{8}\Xi\lambda_2^2-\tfrac{13\ii\beta_4}{32}\lambda_2^4$ \\
$C_4$ & $-\tfrac32\lambda_1+\tfrac{35\beta_3}{16}\lambda_2^3$ & $-\lambda_1+\tfrac{13\beta_3}{12}\lambda_2^3$ & $-\tfrac12\lambda_1+\tfrac{17\beta_3}{48}\lambda_2^3$ \\
$C_5$ & $\tfrac{g_0-\alpha}{2}-\tfrac{5\ii}{4}\Xi\lambda_2^2+\tfrac{163\ii\beta_4}{192}\lambda_2^4$ & $\tfrac{g_0-\alpha}{2}-\tfrac{3\ii}{4}\Xi\lambda_2^2+\tfrac{13\ii\beta_4}{32}\lambda_2^4$ & $\tfrac{g_0-\alpha}{2}-\tfrac{\ii}{4}\Xi\lambda_2^2+\tfrac{23\ii\beta_4}{192}\lambda_2^4$ \\
$C_6$ & $\tfrac12\lambda_1-\tfrac{11\beta_3}{48}\lambda_2^3$ & $\tfrac12\lambda_1-\tfrac{\beta_3}{6}\lambda_2^3$ & $\tfrac12\lambda_1-\tfrac{5\beta_3}{48}\lambda_2^3$ \\
$C_7$ & $-\tfrac{\ii}{8}\Xi\lambda_2^2+\tfrac{5\ii\beta_4}{96}\lambda_2^4$ & $-\tfrac{\ii}{8}\Xi\lambda_2^2+\tfrac{5\ii\beta_4}{96}\lambda_2^4$ & $-\tfrac{\ii}{8}\Xi\lambda_2^2+\tfrac{5\ii\beta_4}{96}\lambda_2^4$ \\
$C_8$ & $\tfrac{\beta_3}{16}\lambda_2^3$ & $\tfrac{\beta_3}{16}\lambda_2^3$ & $\tfrac{\beta_3}{16}\lambda_2^3$ \\
$C_9$ & $-\tfrac{5\ii\beta_4}{128}\lambda_2^4$ & $-\tfrac{5\ii\beta_4}{128}\lambda_2^4$ & $-\tfrac{5\ii\beta_4}{128}\lambda_2^4$ \\
\bottomrule
\end{tabular}
\end{table}

\begin{table}[htbp]
\centering\small
\caption{Coefficients for $n=-1/2$, $m\in\{0,1/2,1,3/2\}$.}
\begin{tabular}{lcccc}
\toprule
Coeff. & $(-1/2,0)$ & $(-1/2,1/2)$ & $(-1/2,1)$ & $(-1/2,3/2)$ \\
\midrule
$C_1$ & $0$ & $-\tfrac{5\ii\beta_4}{128}\lambda_2^4$ & $0$ & $\tfrac{3\ii\beta_4}{128}\lambda_2^4$ \\
$C_2$ & $0$ & $\tfrac{\beta_3}{16}\lambda_2^3$ & $0$ & $-\tfrac{\beta_3}{16}\lambda_2^3$ \\
$C_3$ & $0$ & $-\tfrac{\ii}{8}\Xi\lambda_2^2+\tfrac{5\ii\beta_4}{96}\lambda_2^4$ & $0$ & $\tfrac{3\ii}{8}\Xi\lambda_2^2+\tfrac{3\ii\beta_4}{32}\lambda_2^4$ \\
$C_4$ & $0$ & $\tfrac12\lambda_1+\tfrac{\beta_3}{48}\lambda_2^3$ & $\lambda_1+\tfrac{5\beta_3}{12}\lambda_2^3$ & $\tfrac32\lambda_1+\tfrac{19\beta_3}{16}\lambda_2^3$ \\
$C_5$ & $\tfrac{g_0-\alpha}{2}+\tfrac{\ii}{4}\Xi\lambda_2^2-\tfrac{\ii\beta_4}{96}\lambda_2^4$ & $\tfrac{g_0-\alpha}{2}+\tfrac{3\ii}{4}\Xi\lambda_2^2+\tfrac{\ii\beta_4}{64}\lambda_2^4$ & $\tfrac{g_0-\alpha}{2}+\tfrac{5\ii}{4}\Xi\lambda_2^2+\tfrac{19\ii\beta_4}{96}\lambda_2^4$ & $\tfrac{g_0-\alpha}{2}+\tfrac{7\ii}{4}\Xi\lambda_2^2+\tfrac{103\ii\beta_4}{192}\lambda_2^4$ \\
$C_6$ & $\tfrac12\lambda_1-\tfrac{\beta_3}{24}\lambda_2^3$ & $\tfrac12\lambda_1+\tfrac{\beta_3}{48}\lambda_2^3$ & $\tfrac12\lambda_1+\tfrac{\beta_3}{12}\lambda_2^3$ & $\tfrac12\lambda_1+\tfrac{7\beta_3}{48}\lambda_2^3$ \\
$C_7$ & $-\tfrac{\ii}{8}\Xi\lambda_2^2+\tfrac{5\ii\beta_4}{96}\lambda_2^4$ & $-\tfrac{\ii}{8}\Xi\lambda_2^2+\tfrac{5\ii\beta_4}{96}\lambda_2^4$ & $-\tfrac{\ii}{8}\Xi\lambda_2^2+\tfrac{5\ii\beta_4}{96}\lambda_2^4$ & $-\tfrac{\ii}{8}\Xi\lambda_2^2+\tfrac{5\ii\beta_4}{96}\lambda_2^4$ \\
$C_8$ & $\tfrac{\beta_3}{16}\lambda_2^3$ & $\tfrac{\beta_3}{16}\lambda_2^3$ & $\tfrac{\beta_3}{16}\lambda_2^3$ & $\tfrac{\beta_3}{16}\lambda_2^3$ \\
$C_9$ & $-\tfrac{5\ii\beta_4}{128}\lambda_2^4$ & $-\tfrac{5\ii\beta_4}{128}\lambda_2^4$ & $-\tfrac{5\ii\beta_4}{128}\lambda_2^4$ & $-\tfrac{5\ii\beta_4}{128}\lambda_2^4$ \\
\bottomrule
\end{tabular}
\end{table}

\subsection{$n=0$}

\begin{table}[htbp]
\centering\small
\caption{Coefficients for $n=0$, $m\in\{-3/2,-1,-1/2\}$.}
\begin{tabular}{lccc}
\toprule
Coeff. & $(0,-3/2)$ & $(0,-1)$ & $(0,-1/2)$ \\
\midrule
$C_1$ & $\tfrac{315\ii\beta_4}{128}\lambda_2^4$ & $\ii\beta_4\lambda_2^4$ & $\tfrac{35\ii\beta_4}{128}\lambda_2^4$ \\
$C_2$ & $-\tfrac{35\beta_3}{16}\lambda_2^3$ & $-\beta_3\lambda_2^3$ & $-\tfrac{5\beta_3}{16}\lambda_2^3$ \\
$C_3$ & $\tfrac{15\ii}{8}\Xi\lambda_2^2-\tfrac{85\ii\beta_4}{32}\lambda_2^4$ & $\ii\Xi\lambda_2^2-\tfrac{7\ii\beta_4}{6}\lambda_2^4$ & $\tfrac{3\ii}{8}\Xi\lambda_2^2-\tfrac{11\ii\beta_4}{32}\lambda_2^4$ \\
$C_4$ & $-\tfrac32\lambda_1+\tfrac{13\beta_3}{8}\lambda_2^3$ & $-\lambda_1+\tfrac{5\beta_3}{6}\lambda_2^3$ & $-\tfrac12\lambda_1+\tfrac{7\beta_3}{24}\lambda_2^3$ \\
$C_5$ & $\tfrac{g_0-\alpha}{2}-\tfrac{3\ii}{4}\Xi\lambda_2^2+\tfrac{13\ii\beta_4}{32}\lambda_2^4$ & $\tfrac{g_0-\alpha}{2}-\tfrac{\ii}{2}\Xi\lambda_2^2+\tfrac{5\ii\beta_4}{24}\lambda_2^4$ & $\tfrac{g_0-\alpha}{2}-\tfrac{\ii}{4}\Xi\lambda_2^2+\tfrac{7\ii\beta_4}{96}\lambda_2^4$ \\
$C_6$ & $0$ & $0$ & $0$ \\
$C_7$ & $0$ & $0$ & $0$ \\
$C_8$ & $0$ & $0$ & $0$ \\
$C_9$ & $0$ & $0$ & $0$ \\
\bottomrule
\end{tabular}
\end{table}

\begin{table}[htbp]
\centering\small
\caption{Coefficients for $n=0$, $m\in\{0,1/2,1,3/2\}$.}
\begin{tabular}{lcccc}
\toprule
Coeff. & $(0,0)$ & $(0,1/2)$ & $(0,1)$ & $(0,3/2)$ \\
\midrule
$C_1$ & $0$ & $-\tfrac{5\ii\beta_4}{128}\lambda_2^4$ & $0$ & $\tfrac{3\ii\beta_4}{128}\lambda_2^4$ \\
$C_2$ & $0$ & $\tfrac{\beta_3}{16}\lambda_2^3$ & $0$ & $-\tfrac{\beta_3}{16}\lambda_2^3$ \\
$C_3$ & $0$ & $-\tfrac{\ii}{8}\Xi\lambda_2^2+\tfrac{5\ii\beta_4}{96}\lambda_2^4$ & $0$ & $\tfrac{3\ii}{8}\Xi\lambda_2^2+\tfrac{\ii\beta_4}{32}\lambda_2^4$ \\
$C_4$ & $0$ & $\tfrac12\lambda_1-\tfrac{\beta_3}{24}\lambda_2^3$ & $\lambda_1+\tfrac{\beta_3}{6}\lambda_2^3$ & $\tfrac32\lambda_1+\tfrac{5\beta_3}{8}\lambda_2^3$ \\
$C_5$ & $\tfrac{g_0-\alpha}{2}$ & $\tfrac{g_0-\alpha}{2}+\tfrac{\ii}{4}\Xi\lambda_2^2-\tfrac{\ii\beta_4}{96}\lambda_2^4$ & $\tfrac{g_0-\alpha}{2}+\tfrac{\ii}{2}\Xi\lambda_2^2+\tfrac{\ii\beta_4}{24}\lambda_2^4$ & $\tfrac{g_0-\alpha}{2}+\tfrac{3\ii}{4}\Xi\lambda_2^2+\tfrac{5\ii\beta_4}{32}\lambda_2^4$ \\
$C_6$ & $0$ & $0$ & $0$ & $0$ \\
$C_7$ & $0$ & $0$ & $0$ & $0$ \\
$C_8$ & $0$ & $0$ & $0$ & $0$ \\
$C_9$ & $0$ & $0$ & $0$ & $0$ \\
\bottomrule
\end{tabular}
\end{table}

\subsection{$n=1/2$}

\begin{table}[htbp]
\centering\small
\caption{Coefficients for $n=1/2$, $m\in\{-3/2,-1,-1/2\}$.}
\begin{tabular}{lccc}
\toprule
Coeff. & $(1/2,-3/2)$ & $(1/2,-1)$ & $(1/2,-1/2)$ \\
\midrule
$C_1$ & $\tfrac{315\ii\beta_4}{128}\lambda_2^4$ & $\ii\beta_4\lambda_2^4$ & $\tfrac{35\ii\beta_4}{128}\lambda_2^4$ \\
$C_2$ & $-\tfrac{35\beta_3}{16}\lambda_2^3$ & $-\beta_3\lambda_2^3$ & $-\tfrac{5\beta_3}{16}\lambda_2^3$ \\
$C_3$ & $\tfrac{15\ii}{8}\Xi\lambda_2^2-\tfrac{65\ii\beta_4}{32}\lambda_2^4$ & $\ii\Xi\lambda_2^2-\tfrac{11\ii\beta_4}{12}\lambda_2^4$ & $\tfrac{3\ii}{8}\Xi\lambda_2^2-\tfrac{9\ii\beta_4}{32}\lambda_2^4$ \\
$C_4$ & $-\tfrac32\lambda_1+\tfrac{17\beta_3}{16}\lambda_2^3$ & $-\lambda_1+\tfrac{7\beta_3}{12}\lambda_2^3$ & $-\tfrac12\lambda_1+\tfrac{11\beta_3}{48}\lambda_2^3$ \\
$C_5$ & $\tfrac{g_0-\alpha}{2}-\tfrac{\ii}{4}\Xi\lambda_2^2+\tfrac{23\ii\beta_4}{192}\lambda_2^4$ & $\tfrac{g_0-\alpha}{2}-\tfrac{\ii}{4}\Xi\lambda_2^2+\tfrac{7\ii\beta_4}{96}\lambda_2^4$ & $\tfrac{g_0-\alpha}{2}-\tfrac{\ii}{4}\Xi\lambda_2^2+\tfrac{11\ii\beta_4}{192}\lambda_2^4$ \\
$C_6$ & $-\tfrac12\lambda_1+\tfrac{5\beta_3}{48}\lambda_2^3$ & $-\tfrac12\lambda_1+\tfrac{\beta_3}{6}\lambda_2^3$ & $-\tfrac12\lambda_1+\tfrac{11\beta_3}{48}\lambda_2^3$ \\
$C_7$ & $\tfrac{3\ii}{8}\Xi\lambda_2^2-\tfrac{5\ii\beta_4}{32}\lambda_2^4$ & $\tfrac{3\ii}{8}\Xi\lambda_2^2-\tfrac{7\ii\beta_4}{32}\lambda_2^4$ & $\tfrac{3\ii}{8}\Xi\lambda_2^2-\tfrac{9\ii\beta_4}{32}\lambda_2^4$ \\
$C_8$ & $-\tfrac{5\beta_3}{16}\lambda_2^3$ & $-\tfrac{5\beta_3}{16}\lambda_2^3$ & $-\tfrac{5\beta_3}{16}\lambda_2^3$ \\
$C_9$ & $\tfrac{35\ii\beta_4}{128}\lambda_2^4$ & $\tfrac{35\ii\beta_4}{128}\lambda_2^4$ & $\tfrac{35\ii\beta_4}{128}\lambda_2^4$ \\
\bottomrule
\end{tabular}
\end{table}

\begin{table}[htbp]
\centering\small
\caption{Coefficients for $n=1/2$, $m\in\{0,1/2,1,3/2\}$.}
\begin{tabular}{lcccc}
\toprule
Coeff. & $(1/2,0)$ & $(1/2,1/2)$ & $(1/2,1)$ & $(1/2,3/2)$ \\
\midrule
$C_1$ & $0$ & $-\tfrac{5\ii\beta_4}{128}\lambda_2^4$ & $0$ & $\tfrac{3\ii\beta_4}{128}\lambda_2^4$ \\
$C_2$ & $0$ & $\tfrac{\beta_3}{16}\lambda_2^3$ & $0$ & $-\tfrac{\beta_3}{16}\lambda_2^3$ \\
$C_3$ & $0$ & $-\tfrac{\ii}{8}\Xi\lambda_2^2+\tfrac{5\ii\beta_4}{96}\lambda_2^4$ & $0$ & $\tfrac{3\ii}{8}\Xi\lambda_2^2-\tfrac{\ii\beta_4}{32}\lambda_2^4$ \\
$C_4$ & $0$ & $\tfrac12\lambda_1-\tfrac{5\beta_3}{48}\lambda_2^3$ & $\lambda_1-\tfrac{\beta_3}{12}\lambda_2^3$ & $\tfrac32\lambda_1+\tfrac{\beta_3}{16}\lambda_2^3$ \\
$C_5$ & $\tfrac{g_0-\alpha}{2}-\tfrac{\ii}{4}\Xi\lambda_2^2+\tfrac{7\ii\beta_4}{96}\lambda_2^4$ & $\tfrac{g_0-\alpha}{2}-\tfrac{\ii}{4}\Xi\lambda_2^2+\tfrac{23\ii\beta_4}{192}\lambda_2^4$ & $\tfrac{g_0-\alpha}{2}-\tfrac{\ii}{4}\Xi\lambda_2^2+\tfrac{19\ii\beta_4}{96}\lambda_2^4$ & $\tfrac{g_0-\alpha}{2}-\tfrac{\ii}{4}\Xi\lambda_2^2+\tfrac{59\ii\beta_4}{192}\lambda_2^4$ \\
$C_6$ & $-\tfrac12\lambda_1+\tfrac{7\beta_3}{24}\lambda_2^3$ & $-\tfrac12\lambda_1+\tfrac{17\beta_3}{48}\lambda_2^3$ & $-\tfrac12\lambda_1+\tfrac{5\beta_3}{12}\lambda_2^3$ & $-\tfrac12\lambda_1+\tfrac{23\beta_3}{48}\lambda_2^3$ \\
$C_7$ & $\tfrac{3\ii}{8}\Xi\lambda_2^2-\tfrac{11\ii\beta_4}{32}\lambda_2^4$ & $\tfrac{3\ii}{8}\Xi\lambda_2^2-\tfrac{13\ii\beta_4}{32}\lambda_2^4$ & $\tfrac{3\ii}{8}\Xi\lambda_2^2-\tfrac{15\ii\beta_4}{32}\lambda_2^4$ & $\tfrac{3\ii}{8}\Xi\lambda_2^2-\tfrac{17\ii\beta_4}{32}\lambda_2^4$ \\
$C_8$ & $-\tfrac{5\beta_3}{16}\lambda_2^3$ & $-\tfrac{5\beta_3}{16}\lambda_2^3$ & $-\tfrac{5\beta_3}{16}\lambda_2^3$ & $-\tfrac{5\beta_3}{16}\lambda_2^3$ \\
$C_9$ & $\tfrac{35\ii\beta_4}{128}\lambda_2^4$ & $\tfrac{35\ii\beta_4}{128}\lambda_2^4$ & $\tfrac{35\ii\beta_4}{128}\lambda_2^4$ & $\tfrac{35\ii\beta_4}{128}\lambda_2^4$ \\
\bottomrule
\end{tabular}
\end{table}

\subsection{$n=1$}

\begin{table}[htbp]
\centering\small
\caption{Coefficients for $n=1$, $m\in\{-3/2,-1,-1/2\}$.}
\begin{tabular}{lccc}
\toprule
Coeff. & $(1,-3/2)$ & $(1,-1)$ & $(1,-1/2)$ \\
\midrule
$C_1$ & $\tfrac{315\ii\beta_4}{128}\lambda_2^4$ & $\ii\beta_4\lambda_2^4$ & $\tfrac{35\ii\beta_4}{128}\lambda_2^4$ \\
$C_2$ & $-\tfrac{35\beta_3}{16}\lambda_2^3$ & $-\beta_3\lambda_2^3$ & $-\tfrac{5\beta_3}{16}\lambda_2^3$ \\
$C_3$ & $\tfrac{15\ii}{8}\Xi\lambda_2^2-\tfrac{45\ii\beta_4}{32}\lambda_2^4$ & $\ii\Xi\lambda_2^2-\tfrac{2\ii\beta_4}{3}\lambda_2^4$ & $\tfrac{3\ii}{8}\Xi\lambda_2^2-\tfrac{7\ii\beta_4}{32}\lambda_2^4$ \\
$C_4$ & $-\tfrac32\lambda_1+\tfrac{\beta_3}{2}\lambda_2^3$ & $-\lambda_1+\tfrac{\beta_3}{3}\lambda_2^3$ & $-\tfrac12\lambda_1+\tfrac{\beta_3}{6}\lambda_2^3$ \\
$C_5$ & $\tfrac{g_0-\alpha}{2}+\tfrac{\ii}{4}\Xi\lambda_2^2-\tfrac{\ii\beta_4}{96}\lambda_2^4$ & $\tfrac{g_0-\alpha}{2}$ & $\tfrac{g_0-\alpha}{2}-\tfrac{\ii}{4}\Xi\lambda_2^2+\tfrac{7\ii\beta_4}{96}\lambda_2^4$ \\
$C_6$ & $-\lambda_1+\tfrac{\beta_3}{12}\lambda_2^3$ & $-\lambda_1+\tfrac{\beta_3}{3}\lambda_2^3$ & $-\lambda_1+\tfrac{7\beta_3}{12}\lambda_2^3$ \\
$C_7$ & $\ii\Xi\lambda_2^2-\tfrac{5\ii\beta_4}{12}\lambda_2^4$ & $\ii\Xi\lambda_2^2-\tfrac{2\ii\beta_4}{3}\lambda_2^4$ & $\ii\Xi\lambda_2^2-\tfrac{11\ii\beta_4}{12}\lambda_2^4$ \\
$C_8$ & $-\beta_3\lambda_2^3$ & $-\beta_3\lambda_2^3$ & $-\beta_3\lambda_2^3$ \\
$C_9$ & $\ii\beta_4\lambda_2^4$ & $\ii\beta_4\lambda_2^4$ & $\ii\beta_4\lambda_2^4$ \\
\bottomrule
\end{tabular}
\end{table}

\begin{table}[htbp]
\centering\small
\caption{Coefficients for $n=1$, $m\in\{0,1/2,1,3/2\}$.}
\label{tab:n1-0113}
\begin{tabular}{lcccc}
\toprule
Coeff. & $(1,0)$ & $(1,1/2)$ & $(1,1)$ & $(1,3/2)$ \\
\midrule
$C_1$ & $0$ & $-\tfrac{5\ii\beta_4}{128}\lambda_2^4$ & $0$ & $\tfrac{3\ii\beta_4}{128}\lambda_2^4$ \\
$C_2$ & $0$ & $\tfrac{\beta_3}{16}\lambda_2^3$ & $0$ & $-\tfrac{\beta_3}{16}\lambda_2^3$ \\
$C_3$ & $0$ & $-\tfrac{\ii}{8}\Xi\lambda_2^2+\tfrac{5\ii\beta_4}{96}\lambda_2^4$ & $0$ & $\tfrac{3\ii}{8}\Xi\lambda_2^2-\tfrac{3\ii\beta_4}{32}\lambda_2^4$ \\
$C_4$ & $0$ & $\tfrac12\lambda_1-\tfrac{\beta_3}{6}\lambda_2^3$ & $\lambda_1-\tfrac{\beta_3}{3}\lambda_2^3$ & $\tfrac32\lambda_1-\tfrac{\beta_3}{2}\lambda_2^3$ \\
$C_5$ & $\tfrac{g_0-\alpha}{2}-\tfrac{\ii}{2}\Xi\lambda_2^2+\tfrac{5\ii\beta_4}{24}\lambda_2^4$ & $\tfrac{g_0-\alpha}{2}-\tfrac{3\ii}{4}\Xi\lambda_2^2+\tfrac{13\ii\beta_4}{32}\lambda_2^4$ & $\tfrac{g_0-\alpha}{2}-\ii\Xi\lambda_2^2+\tfrac{2\ii\beta_4}{3}\lambda_2^4$ & $\tfrac{g_0-\alpha}{2}-\tfrac{5\ii}{4}\Xi\lambda_2^2+\tfrac{95\ii\beta_4}{96}\lambda_2^4$ \\
$C_6$ & $-\lambda_1+\tfrac{5\beta_3}{6}\lambda_2^3$ & $-\lambda_1+\tfrac{13\beta_3}{12}\lambda_2^3$ & $-\lambda_1+\tfrac{4\beta_3}{3}\lambda_2^3$ & $-\lambda_1+\tfrac{19\beta_3}{12}\lambda_2^3$ \\
$C_7$ & $\ii\Xi\lambda_2^2-\tfrac{7\ii\beta_4}{6}\lambda_2^4$ & $\ii\Xi\lambda_2^2-\tfrac{17\ii\beta_4}{12}\lambda_2^4$ & $\ii\Xi\lambda_2^2-\tfrac{5\ii\beta_4}{3}\lambda_2^4$ & $\ii\Xi\lambda_2^2-\tfrac{23\ii\beta_4}{12}\lambda_2^4$ \\
$C_8$ & $-\beta_3\lambda_2^3$ & $-\beta_3\lambda_2^3$ & $-\beta_3\lambda_2^3$ & $-\beta_3\lambda_2^3$ \\
$C_9$ & $\ii\beta_4\lambda_2^4$ & $\ii\beta_4\lambda_2^4$ & $\ii\beta_4\lambda_2^4$ & $\ii\beta_4\lambda_2^4$ \\
\bottomrule
\end{tabular}
\end{table}

\subsection{$n=3/2$}

\begin{table}[htbp]
\centering\small
\caption{Coefficients for $n=3/2$, $m\in\{-3/2,-1,-1/2\}$.}
\begin{tabular}{lccc}
\toprule
Coeff. & $(3/2,-3/2)$ & $(3/2,-1)$ & $(3/2,-1/2)$ \\
\midrule
$C_1$ & $\tfrac{315\ii\beta_4}{128}\lambda_2^4$ & $\ii\beta_4\lambda_2^4$ & $\tfrac{35\ii\beta_4}{128}\lambda_2^4$ \\
$C_2$ & $-\tfrac{35\beta_3}{16}\lambda_2^3$ & $-\beta_3\lambda_2^3$ & $-\tfrac{5\beta_3}{16}\lambda_2^3$ \\
$C_3$ & $\tfrac{15\ii}{8}\Xi\lambda_2^2-\tfrac{25\ii\beta_4}{32}\lambda_2^4$ & $\ii\Xi\lambda_2^2-\tfrac{5\ii\beta_4}{12}\lambda_2^4$ & $\tfrac{3\ii}{8}\Xi\lambda_2^2-\tfrac{5\ii\beta_4}{32}\lambda_2^4$ \\
$C_4$ & $-\tfrac32\lambda_1-\tfrac{\beta_3}{16}\lambda_2^3$ & $-\lambda_1+\tfrac{\beta_3}{12}\lambda_2^3$ & $-\tfrac12\lambda_1+\tfrac{5\beta_3}{48}\lambda_2^3$ \\
$C_5$ & $\tfrac{g_0-\alpha}{2}+\tfrac{3\ii}{4}\Xi\lambda_2^2+\tfrac{\ii\beta_4}{64}\lambda_2^4$ & $\tfrac{g_0-\alpha}{2}+\tfrac{\ii}{4}\Xi\lambda_2^2-\tfrac{\ii\beta_4}{96}\lambda_2^4$ & $\tfrac{g_0-\alpha}{2}-\tfrac{\ii}{4}\Xi\lambda_2^2+\tfrac{23\ii\beta_4}{192}\lambda_2^4$ \\
$C_6$ & $-\tfrac32\lambda_1-\tfrac{\beta_3}{16}\lambda_2^3$ & $-\tfrac32\lambda_1+\tfrac{\beta_3}{2}\lambda_2^3$ & $-\tfrac32\lambda_1+\tfrac{17\beta_3}{16}\lambda_2^3$ \\
$C_7$ & $\tfrac{15\ii}{8}\Xi\lambda_2^2-\tfrac{25\ii\beta_4}{32}\lambda_2^4$ & $\tfrac{15\ii}{8}\Xi\lambda_2^2-\tfrac{45\ii\beta_4}{32}\lambda_2^4$ & $\tfrac{15\ii}{8}\Xi\lambda_2^2-\tfrac{65\ii\beta_4}{32}\lambda_2^4$ \\
$C_8$ & $-\tfrac{35\beta_3}{16}\lambda_2^3$ & $-\tfrac{35\beta_3}{16}\lambda_2^3$ & $-\tfrac{35\beta_3}{16}\lambda_2^3$ \\
$C_9$ & $\tfrac{315\ii\beta_4}{128}\lambda_2^4$ & $\tfrac{315\ii\beta_4}{128}\lambda_2^4$ & $\tfrac{315\ii\beta_4}{128}\lambda_2^4$ \\
\bottomrule
\end{tabular}
\end{table}

\begin{table}[htbp]
\centering\small
\caption{Coefficients for $n=3/2$, $m\in\{0,1/2,1,3/2\}$.}
\label{tab:n32-3-32}
\begin{tabular}{lcccc}
\toprule
Coeff. & $(3/2,0)$ & $(3/2,1/2)$ & $(3/2,1)$ & $(3/2,3/2)$ \\
\midrule
$C_1$ & $0$ & $-\tfrac{5\ii\beta_4}{128}\lambda_2^4$ & $0$ & $\tfrac{3\ii\beta_4}{128}\lambda_2^4$ \\
$C_2$ & $0$ & $\tfrac{\beta_3}{16}\lambda_2^3$ & $0$ & $-\tfrac{\beta_3}{16}\lambda_2^3$ \\
$C_3$ & $0$ & $-\tfrac{\ii}{8}\Xi\lambda_2^2+\tfrac{5\ii\beta_4}{96}\lambda_2^4$ & $0$ & $\tfrac{3\ii}{8}\Xi\lambda_2^2-\tfrac{5\ii\beta_4}{32}\lambda_2^4$ \\
$C_4$ & $0$ & $\tfrac12\lambda_1-\tfrac{11\beta_3}{48}\lambda_2^3$ & $\lambda_1-\tfrac{7\beta_3}{12}\lambda_2^3$ & $\tfrac32\lambda_1-\tfrac{17\beta_3}{16}\lambda_2^3$ \\
$C_5$ & $\tfrac{g_0-\alpha}{2}-\tfrac{3\ii}{4}\Xi\lambda_2^2+\tfrac{13\ii\beta_4}{32}\lambda_2^4$ & $\tfrac{g_0-\alpha}{2}-\tfrac{5\ii}{4}\Xi\lambda_2^2+\tfrac{163\ii\beta_4}{192}\lambda_2^4$ & $\tfrac{g_0-\alpha}{2}-\tfrac{7\ii}{4}\Xi\lambda_2^2+\tfrac{139\ii\beta_4}{96}\lambda_2^4$ & $\tfrac{g_0-\alpha}{2}-\tfrac{9\ii}{4}\Xi\lambda_2^2+\tfrac{141\ii\beta_4}{64}\lambda_2^4$ \\
$C_6$ & $-\tfrac32\lambda_1+\tfrac{13\beta_3}{8}\lambda_2^3$ & $-\tfrac32\lambda_1+\tfrac{35\beta_3}{16}\lambda_2^3$ & $-\tfrac32\lambda_1+\tfrac{11\beta_3}{4}\lambda_2^3$ & $-\tfrac32\lambda_1+\tfrac{53\beta_3}{16}\lambda_2^3$ \\
$C_7$ & $\tfrac{15\ii}{8}\Xi\lambda_2^2-\tfrac{85\ii\beta_4}{32}\lambda_2^4$ & $\tfrac{15\ii}{8}\Xi\lambda_2^2-\tfrac{105\ii\beta_4}{32}\lambda_2^4$ & $\tfrac{15\ii}{8}\Xi\lambda_2^2-\tfrac{125\ii\beta_4}{32}\lambda_2^4$ & $\tfrac{15\ii}{8}\Xi\lambda_2^2-\tfrac{145\ii\beta_4}{32}\lambda_2^4$ \\
$C_8$ & $-\tfrac{35\beta_3}{16}\lambda_2^3$ & $-\tfrac{35\beta_3}{16}\lambda_2^3$ & $-\tfrac{35\beta_3}{16}\lambda_2^3$ & $-\tfrac{35\beta_3}{16}\lambda_2^3$ \\
$C_9$ & $\tfrac{315\ii\beta_4}{128}\lambda_2^4$ & $\tfrac{315\ii\beta_4}{128}\lambda_2^4$ & $\tfrac{315\ii\beta_4}{128}\lambda_2^4$ & $\tfrac{315\ii\beta_4}{128}\lambda_2^4$ \\
\bottomrule
\end{tabular}
\end{table}

%======================================================================
\section*{Data and Code Availability}
%======================================================================
The coefficient expressions in Appendix~\ref{app:tables} were obtained by direct exact
substitution of $(n,m)$ into the closed-form expressions of Section~5 and are fully
reproducible from that section alone. No numerical PINN experiments were executed for
this paper; Section~\ref{sec:pinn} specifies the methodology for future numerical work.

%======================================================================
\bibliographystyle{unsrtnat}
\begin{thebibliography}{99}

\bibitem[Raissi et al.(2019)]{raissi2019}
M.~Raissi, P.~Perdikaris, G.~E.~Karniadakis,
\emph{Physics-informed neural networks: A deep learning framework for solving forward and
inverse problems involving nonlinear partial differential equations},
Journal of Computational Physics \textbf{378} (2019) 686--707.

\bibitem[Zhou and Yan(2021)]{zhou2021}
Z.~Zhou, Z.~Yan,
\emph{Deep learning neural networks for the third-order nonlinear Schr\"odinger equation:
solitons, breathers, and rogue waves},
arXiv:2104.14809 (2021).

\bibitem[Fang et al.(2021)]{fang2021}
Y.~Fang, G.-Z.~Wu, Y.-Y.~Wang, C.-Q.~Dai,
\emph{Data-driven femtosecond optical soliton excitations and parameters discovery of the
high-order NLSE using the PINN},
arXiv:2103.16297 (2021).

\bibitem[Physics-Informed Neural Networks Review(2025)]{pinnreview2025}
\emph{Physics-Informed Neural Networks for Higher-Order Nonlinear Schr\"odinger Equations:
Soliton Dynamics in External Potentials}, preprint (2025).

\bibitem[GRAVIS(2026)]{gravis2026}
\emph{Physics-informed neural networks for nonlinear Schr\"odinger localized wave problems
with gradient-residual adaptive collocation sampling},
Nonlinear Dynamics (2026).

\bibitem[Baldwin et al.(2004)]{baldwin2004}
D.~Baldwin, \"U.~G\"okta\c{s}, W.~Hereman, L.~Hong, R.~S.~Martino, J.~C.~Miller,
\emph{Symbolic computation of exact solutions expressible in hyperbolic and elliptic
functions for nonlinear PDEs},
Journal of Symbolic Computation \textbf{37}(6) (2004) 669--705.

\bibitem[Hereman and Takaoka(1990)]{hereman1990}
W.~Hereman, M.~Takaoka,
\emph{Solitary wave solutions of nonlinear evolution and wave equations using a direct
algebraic method},
Journal of Physics A: Mathematical and General \textbf{23} (1990) 4805--4822.

\end{thebibliography}

\end{document}